\section{研究動機與背景}
近年來，基於深度學習的文字轉語音 (Text-to-Speech, TTS) 與語音轉換 (Voice Conversion, VC) 技術取得了突破性的發展，所生成的合成語音在自然度與音色相似度上已難以與真實人類語音區分。這項技術的進步雖然為語音助理、有聲書等應用帶來便利，卻也同時催生了嚴重的「深偽語音 (Audio Deepfake)」威脅。不法分子頻繁將此技術應用於電話詐騙、身份盜用或虛假訊息傳播，對自動說話人驗證 (Automatic Speaker Verification, ASV) 系統與社會信任構成了嚴峻的安全挑戰。

為應對此一威脅，學界與業界積極投入語音深偽偵測技術的研發。然而，現存的多數偵測系統在面對訓練集從未見過的攻擊演算法，或是真實世界中存在各種背景噪音與通訊壓縮失真的場景時，往往面臨嚴重的效能衰退。這種「跨域泛化能力 (Cross-domain Generalization)」的匱乏，已成為當前深偽偵測防禦系統從實驗室走向實務部署的最大瓶頸。

\section{研究目的與核心問題}
隨著語音基礎模型 (Speech Foundation Models) 的興起，深偽偵測領域逐漸形成兩種具代表性的訓練範式 (Training Paradigms)：
第一種為目前廣泛採用的「預訓練—微調（Pre-train, Fine-tune）」範式，此類模型（如 W2V-AASIST 與 MFA-Conformer）先在大規模語音資料上進行自監督預訓練，再針對深偽偵測任務進行全參數或部分微調，以連續的波形或頻譜特徵捕捉細微的偽造痕跡；
第二種為新興的「Prompt Tuning」範式，此類方法（如 SpeechPrompt v2）凍結預訓練語音語言模型的主體參數，僅透過極少量可學習的 Prompt Token 引導模型完成分類任務，具備極低的訓練參數量。由於現行語音語言模型採離散 Token 作為輸入，Prompt Tuning 範式須先對聲學特徵進行量化離散化，而此過程所帶來的聲學細節損失，是制約其跨域偵測能力的關鍵因素。

本研究的核心目的，在於以跨域泛化能力為核心評估標準，對此兩種訓練範式進行系統性的比較與架構診斷。我們試圖解答以下核心問題：
(1) 在面對未知的跨資料集攻擊時，預訓練—微調與 Prompt Tuning 範式各自呈現何種防禦特性與先天極限？
(2) Prompt Tuning 架構的跨域表現落差，究竟源於後端語言模型的推論能力不足，抑或前端離散量化所造成的聲學資訊流失？
(3) 對於預訓練—微調架構，單純擴增模型規模是否能有效解決跨域效能飽和的問題？

\section{研究範圍與貢獻}
為客觀解答上述問題，本研究選用具代表性的三種架構：W2V-AASIST、MFA-Conformer 與 SpeechPrompt v2，並橫跨 ASVspoof 2019/2021 (包含 Logical Access 與 DeepFake)、真實世界資料集 InTheWild，以及大規模多語言資料集 MLAAD，進行了廣泛且嚴謹的跨域交叉測試。

本研究之主要貢獻總結如下：
\begin{enumerate}
    \item \textbf{建立雙範式的系統性對比視角}：首次在統一的評估基準下（相同的運算精度、損失函數與亂數種子嚴格平均），全面比較預訓練—微調與 Prompt Tuning 兩大訓練範式在深偽偵測任務上的優劣與跨域泛化特性。
    \item \textbf{實證離散特徵架構之效能瓶頸}：透過設計 Prompt Tuning 與 Full Fine-Tuning 對比，以及提出「軟性離散特徵」與「可學習群集中心」等消融實驗，成功實證 SpeechPrompt 等架構的效能天花板係源自前端離散化過程中的聲學細節損失，而非後端語言模型的學習限制。
    \item \textbf{剖析預訓練—微調架構之跨域瓶頸}：透過系統性的規模擴容與後端架構消融實驗，揭示預訓練—微調範式下跨域效能的瓶頸並非源於前端模型規模的不足，而在於後端特徵融合架構的設計能力。
\end{enumerate}

\section{報告結構}
本報告後續章節編排如下：
第二章回顧主流的語音合成技術，以解析深偽攻擊之發展脈絡。
第三章回顧語音深偽偵測技術的發展脈絡，分別介紹預訓練—微調範式下的代表性模型（W2V-AASIST、MFA-Conformer），以及 Prompting 範式下的核心架構（GSLM 與 SpeechPrompt v2）。
第四章詳述實驗設置、資料集配置，並透過主實驗與消融實驗深入剖析兩大訓練範式在跨域偵測上的表現與瓶頸。
第五章展示深偽偵測模型之實務前端應用系統設計。
第六章總結本研究之各項實驗發現，並進行綜合討論。
第七章基於本研究成果，提出對未來語音深偽防禦技術發展之具體建議與展望。
